% 2010 年天津卷文科第 8 题：区间 [−π/6, 5π/6] 上的 y=A sin(ωx+φ) 图象。
% 来源 PDF：普通高考/2010/2010天津文.pdf（上游勘误版），SHA-256 e6b06813e17b7c332b651b8784c8fab0fcad0c8666ea6032a2958917c55307c9。
% 原页：第 1 页（A3 四栏）第 2 栏上部，书页码 637。
% 重绘依据：曲线过 (−π/6,0)、(π/3,0)、(5π/6,0)，在 π/12 取最大值 1，在 7π/12 取最小值 −1。
\begin{tikzpicture}[semithick, x=2.4cm, y=1.5cm, every node/.style={inner sep=1.2pt, font=\small}]
  \draw[-stealth] (-0.62,0) -- (3.0,0) node[below] {$x$};
  \draw[-stealth] (0,-1.35) -- (0,1.45) node[left] {$y$};
  \node[below left] at (0,0) {$O$};
  \draw[thick,domain=-0.5236:2.6180,samples=160] plot (\x,{sin(2*\x r + pi/3 r)});
  \draw[dashed] (0,-1) -- (1.5,-1);
  \draw (-0.5236,0.04) -- (-0.5236,-0.04) node[below left] {$-\dfrac\pi6$};
  \draw (1.0472,0.04) -- (1.0472,-0.04) node[below] {$\dfrac\pi3$};
  \draw (2.6180,0.04) -- (2.6180,-0.04) node[below] {$\dfrac{5\pi}6$};
  \draw (-0.04,1) -- (0.04,1) node[left] {$1$};
  \draw (-0.04,-1) -- (0.04,-1) node[left] {$-1$};
\end{tikzpicture}
